% Part: turing-machines
% Chapter: machines-computations
% Section: halting-states

\documentclass[../../../include/open-logic-section]{subfiles}

\begin{document}

\olfileid{tur}{mac}{hal}
\olsection{થંભવાની અવસ્થાઓ}

\begin{explain}
આપણે મશીનને એવી રીતે વ્યાખ્યાયિત કર્યું છે કે અમલમાં મૂકવા
કોઈ સૂચના ન હોય ત્યારે જ તે થંભે. જોકે ટ્યુરિંગ મશીનનાં
સામાન્ય નિરૂપણોમાં એક ખાસ \emph{થંભવાની અવસ્થા}~$h$
હોય છે, જેમાં $h \in Q$.

થંભવાની અવસ્થા પાછળનો વિચાર સરળ છે: મશીનનું કામ પૂરું
થાય (તે ઇનપુટ સ્વીકારવા તૈયાર થાય અથવા આઉટપુટ
લખવાનું પૂરું કરે), ત્યારે તે અવસ્થા~$h$માં જઈને થંભે.
કેટલાંક મશીનોમાં થંભવાની બે અવસ્થાઓ હોય છે: એક
ઇનપુટ સ્વીકારવા માટે અને બીજી ઇનપુટ નકારવા માટે.
\end{explain}

\begin{ex}\emph{થંભવાની અવસ્થાઓ}.
આ વિચાર સ્પષ્ટ કરવા પહેલાં સમ મશીનમાં ફેરફાર કરીએ.
ઇનપુટ સમ હોય ત્યારે મશીન અવસ્થા~$q_0$માં થંભે તે
બદલે, તેને થંભવાની અવસ્થામાં મોકલતી સૂચના ઉમેરી
શકીએ છીએ.
\[
\begin{tikzpicture}[->,>=stealth',shorten >=1pt,auto,node distance=2.8cm,
                    semithick]
  \tikzstyle{every state}=[fill=none,draw=black,text=black]

  \node[initial, state]         (A)                     {$q_0$};
  \node[state]         (B) [right of=A] {$q_1$};
  \node[state]         (C) [below of=A] {$h$};

  \path (A) edge [bend left] node {\TMtrans{\TMstroke}{\TMstroke}{R}} (B)
      edge node {\TMtrans{\TMblank}{\TMblank}{N}} (C)
        (B) edge [loop above] node {\TMtrans{\TMblank}{\TMblank}{R}} (B)
            edge [bend left] node {\TMtrans{\TMstroke}{\TMstroke}{R}} (A);
\end{tikzpicture}
\]

હવે ઉદાહરણને વધુ વિસ્તૃત કરીએ. ઇનપુટ વિષમ છે એવું
મશીન નક્કી કરે ત્યારે તે ક્યારેય થંભતું નથી. આ
ચક્રમાં ચાલવાની સૂચનાને નકારવાની અવસ્થા~$r$માં જવાની
સૂચનાથી બદલીને મશીનમાં \emph{નકારવાની} અવસ્થા
ઉમેરી શકીએ છીએ.
\[
\begin{tikzpicture}[->,>=stealth',shorten >=1pt,auto,node distance=2.8cm,
                    semithick]
  \tikzstyle{every state}=[fill=none,draw=black,text=black]

  \node[initial,state]         (A)                     {$q_0$};
  \node[state]         (B) [right of=A] {$q_1$};
  \node[state]         (C) [below of=A] {$h$};
  \node[state]         (D) [below of=B] {$r$};

  \path (A) edge [bend left] node {\TMtrans{\TMstroke}{\TMstroke}{R}} (B)
      edge node {\TMtrans{\TMblank}{\TMblank}{N}} (C)
        (B) edge node {\TMtrans{\TMblank}{\TMblank}{N}} (D)
            edge [bend left] node {\TMtrans{\TMstroke}{\TMstroke}{R}} (A);
\end{tikzpicture}
\]
\end{ex}

\begin{explain}
આવા કિસ્સામાં ખાસ થંભવાની અવસ્થા ઉમેરવી લાભકારક છે:
મશીન કયા ઇનપુટ સ્વીકારે કે નકારે છે તે સ્પષ્ટ થાય છે.
પરંતુ ખાસ થંભવાની અવસ્થા ઉમેરવાથી મશીનની સંગણનશક્તિ
વધતી નથી એ નોંધવું જરૂરી છે. એવી જ રીતે, થંભવાનો
ઓછો ઔપચારિક ખ્યાલ પણ પોતાના લાભ ધરાવે છે. આ
પ્રકરણમાં અત્યાર સુધી વપરાયેલી થંભવાની વ્યાખ્યા
\emph{અટકવાની સમસ્યા}ની સાબિતી સહજ અને રજૂ કરવામાં
સરળ બનાવે છે. તેથી આપણે આપણી મૂળ વ્યાખ્યા જ
ચાલુ રાખીએ છીએ.
\end{explain}

\end{document}
